% ============================================================
%  CONCLUSION GÉNÉRALE
% ============================================================
\chapter*{Conclusion g\'en\'erale}
\addcontentsline{toc}{chapter}{Conclusion g\'en\'erale}
\markboth{Conclusion g\'en\'erale}{}
\vspace{0.3cm}

Ce Projet de Fin d'\'Etudes, r\'ealis\'e pour le compte de Capgemini au sein de Capgemini Tunisie, marque l'aboutissement de mon cursus d'ing\'enieur \`a Tek-Up et repr\'esente une mise en pratique concr\`ete des comp\'etences techniques et m\'ethodologiques acquises durant ma formation.

Inscrit au c\oe ur de la modernisation de la gestion des projets a\'eronautiques dans un contexte de conformit\'e \`a la norme DO-178C, ce travail visait \`a optimiser un processus de suivi et de r\'evision de code historiquement manuel et complexe. La plateforme SmartPM r\'epond \`a cet enjeu en rationalisant l'ensemble de la cha\^ine op\'erationnelle et en apportant un cadre unifi\'e, automatis\'e et collaboratif pour la coordination entre les d\'eveloppeurs, les relecteurs et les managers.

Le projet a \'et\'e men\'e selon une approche agile Scrum, structur\'ee autour de deux releases compl\'ementaires :
\begin{itemize}
  \item Release 1 : mise en place du socle fonctionnel (authentification s\'ecuris\'ee, gestion de la hi\'erarchie des projets, cr\'eation et suivi des t\^aches et des revues techniques) ;
  \item Release 2 : automatisation de l'assistance contextuelle par l'IA (RAG multi-LLM) et introduction d'une infrastructure DevOps compl\`ete (Kubernetes, Monitoring) pour le pilotage op\'erationnel.
\end{itemize}

Cette d\'emarche a abouti \`a des r\'esultats tangibles : diminution des t\^aches manuelles, meilleure synchronisation entre acteurs, fiabilisation du suivi des revues techniques et mise \`a disposition d'indicateurs de performance permettant un pilotage \'eclair\'e du processus. La plateforme offre d\'esormais une base technique robuste, con\c{c}ue pour \'evoluer et s'adapter \`a de futurs besoins.

Au-del\`a des aspects techniques, ce projet a \'et\'e une exp\'erience humaine enrichissante, m\^elant collaboration, compr\'ehension fine des enjeux m\'etiers et rigueur dans la conception logicielle. Il m'a permis de consolider mes comp\'etences en architecture logicielle, en int\'egration de syst\`emes et en pratiques DevOps, tout en d\'ecouvrant les r\'ealit\'es op\'erationnelles d'un grand groupe.

Les perspectives d'\'evolution de SmartPM sont nombreuses : \'elargissement \`a de nouveaux cas d'usage, int\'egration d'algorithmes d'automatisation plus avanc\'es, am\'elioration de l'exp\'erience utilisateur et exploration de nouveaux indicateurs de pilotage.

Au final, SmartPM constitue aujourd'hui un socle solide et \'evolutif qui modernise un maillon essentiel de l'ing\'enierie logicielle a\'eronautique. Ce projet contribue directement \`a l'am\'elioration de la performance op\'erationnelle et, \`a terme, \`a une meilleure garantie de s\'ecurit\'e et de conformit\'e pour les clients de Capgemini.
